%% ============================================================
%% 分片 30：原文第 10 节「The Bogomolny Equations And The Space Of Hecke Modifications」（原文 glshort.tex 起于第 8120 行）
%% 由 tools/harvmac2tex.py 从 arXiv 源码（hep-th/0604151v3）自动生成。
%% 只译英文散文；公式、编号、\label、\cite 一律不改。
%% ============================================================

\subsection{The Extended Bogomolny Equations}

\label{sec:extbog}

So far, in our study of the action of 't Hooft operators on
branes,  we have considered only supersymmetric solutions that are
time-independent and time-reversal invariant.  With this
restriction, the BPS equations \eqref{eq:oops}\ reduce to the ordinary
Bogomolny equations, and the 't Hooft operators reduce to the
usual geometric Hecke operators.

The assumption of time-independence can actually be justified in
general, but requiring time-reversal invariance puts a severe
restriction on the allowed boundary conditions.  It is only
adequate if branes are chosen (as we did in section 第~\ref{sec:opprod}~节\ in
discussing the product of 't Hooft operators) to ensure that the
relevant solutions have $\varphi=0$. The most important $A$-branes
for the geometric Langlands duality (the generic fibers of the
Hitchin fibrations and the c.c. brane described in section
第~\ref{sec:abranes}~节) do not have their support limited to $\varphi=0$, so in
order to study the action of 't Hooft operators on them, we need
to drop the condition of time-reversal invariance.  This is the
purpose of the following discussion.  Considerably more detail
will appear elsewhere \cite{ref26}.

We consider again the BPS equations on a four-manifold
$M=\RR\times W$ and write the Higgs field $\phi$ as $\phi_0\, dx^0
+\pi^*\tilde\phi$, where $\pi$ is the projection to $W$, and
$\tilde\phi$ is a one-form Higgs field on $W$. Likewise, we write
the gauge field as $A=A_0\,dx^0+\tilde A$, where $\tilde A$ is a
three-dimensional connection with curvature $\tilde F$. The
time-independent BPS equations for $t=1$ read
% 式 (10.35)
\setcounter{equation}{34}
\begin{equation}\label{eq:BPSstaticgen}
\begin{aligned}\tilde F-\tilde\phi\wedge\tilde\phi&=
\star\left(D\phi_0-[A_0,\tilde\phi]\right),\\ \star
D\tilde\phi&=[\phi_0,\tilde\phi]+D A_0,\\
D^*\tilde\phi+[A_0,\phi_0]&=0.\end{aligned}
\end{equation}
 The exterior derivative $D$, the
Hodge star operator $\star$, and the operator $D^*=\star D\star$
are all understood in the three-dimensional sense.

These equations can be greatly simplified using suitable vanishing
theorems.  In fact, one can show that $\phi_1=A_0=0$ using the
same arguments, and the same type of assumptions, as in the
time-reversal invariant case (recall eqn. \eqref{eq:wungo}).

So we set $\phi_1=A_0=0$, choose a local complex coordinate
$z=x^2+ix^3$ on $C$, denote $y=x^1$,  and work in the gauge
$\tilde A_y=0$. We also now drop the tildes from our notation.
And, as in section 第~\ref{sec:theck}~节, we take the metric on $W$ to be $ds^2=h
|dz|^2+ dy^2$, for some function $h$. Then the extended Bogomolny
equations \eqref{eq:BPSstaticgen}\ reduce to
% 式 (10.36)
\setcounter{equation}{35}
\begin{equation}\label{eq:BPSextsimple}
\begin{aligned}D_\bz\phi_z&=0,\\
\partial_y A_\bz&=-i D_\bz \phi_0,\\
\partial_y\phi_z&=-i[\phi_z,\phi_0],\\
F_{z\bz}-[\phi_z,\phi_\bz]&={ih\over 2}\partial_y\phi_0.\end{aligned}
\end{equation}

The first of the equations \eqref{eq:BPSextsimple}\ says that
$\varphi=\phi_z\,dz$, restricted to $C_y=\{y\}\times C$, is a
holomorphic section of $\End(E)\otimes K_C$ for any $y\in I$. So
denoting as $E$ the holomorphic bundle over $C_y$ defined by the
$\bar\partial$ operator $\bar D=d\bar z(\partial_{\bar z}+A_{\bar
z})$, the pair $(E,\varphi)$ is a Higgs bundle for any $y$. The
next two equations have an interpretation that should be easy to
anticipate based on our experience with the ordinary Bogomolny
equations in section 第~\ref{sec:theck}~节: they say that evolution in the $y$
direction changes the Higgs pair $(E,\varphi)$ only by a gauge
transformation.  Thus until we encounter the singularity
associated with an 't Hooft operator, the point in $\MH$
determined by a solution of the extended Bogomolny equations is
independent of $y$.

Now suppose that there is an 't Hooft operator of weight
$^L\negthinspace w$ at, say, $y=y_0$, $z=z_0$.  When $y$ is
increased past $y=y_0$, the bundle $E$ undergoes an ordinary Hecke
modification, as described in section 第~\ref{sec:theck}~节.  In fact, the
evolution equation for $A_{\bar z}$ is the same as it was in
section 第~\ref{sec:theck}~节, and so is the singularity due to the 't Hooft
operator. The pair $(E,\varphi)$ likewise undergoes what we will
call a Hecke modification.  The role of $\varphi$ is, however,
quite different from the role of $E$.

A holomorphic bundle can be modified at a single point, while
preserving holomorphy, and this is what the Hecke or 't Hooft
operators do.  There is a moduli space of possible such
modifications, the familiar space $\bar {\EUY}({}^L\negthinspace
w;z_0)$ parameterizing singular solutions of the Bogomolny
equations (and isomorphic to a Schubert variety). But a
holomorphic field, in this case a holomorphic section $\varphi$ of
${\rm ad}(E)\otimes K_C$, cannot be modified at a point while
preserving holomorphy. Hence, if $\varphi$ is given away from
$z=0$, it has at most a unique holomorphic extension over $z=0$.

So including $\varphi$ does not cause the Hecke modification due
to the 't Hooft operator to depend on additional parameters.  On
the contrary, including $\varphi$ generically eliminates some
parameters.  The reason is that even if $\varphi$ is holomorphic
at $z=z_0$ for $y<y_0$, it will generically have a pole at $z=z_0$
for $y>y_0$.  To see how this happens, let us take $G=U(N)$ and
consider an 't Hooft operator with weight $^L\negthinspace
w=(m_1,\ldots,m_N)$. The holomorphic type of the bundle $E$ jumps
in crossing $y=y_0$; we write $E_-$ and $E_+$ for the bundles
``before'' and ``after.'' Relative to some local decomposition
$E_-\cong \oplus_{i=1}^N \cL_i,$ the bundle $E_+$ has a local
decomposition as $\oplus_{i=1}^N\cL_i\otimes \CO(z_0)^{m_i}$. That
is, if $\cL_i$ is locally generated by a holomorphic section
$s_i$, $i=1,\ldots,N$, then $E_+$ is generated by holomorphic
sections % 式 (10.37)
\setcounter{equation}{36}
\begin{equation}\label{eq:zilfog}
t_i=(z-z_0)^{-m_i} s_i.
\end{equation}
 Here we used a
natural isomorphism between $E_+$ and $E_-$ for $z\neq 0$.

Suppose that $\varphi$ is holomorphic at $z_0$ for $y<y_0$. The
products $\varphi \,s_i$ are holomorphic sections of $E_- \otimes
K_C$ (where $K_C$ will play no essential role as we can trivialize
it near $z=z_0$), so they can be expanded in terms of the $s_i$:
% 式 (10.38)
\setcounter{equation}{37}
\begin{equation}\label{eq:bongo}
\varphi\cdot s_i=\sum_j f_i{}^j\,s_j.
\end{equation}
  Holomorphy of
$\varphi$ at $z=z_0$ for $y<y_0$ says that the $f_i{}^j$ (which
are sections of $K_C$) are holomorphic at $z=z_0$.  They are
subject to no other constraint.

Similarly, for $y>y_0$, we can expand $\varphi\,t_i$ as a linear
combination of the $t_j$.  Given the relation \eqref{eq:zilfog}\ between the
$s_i$ and the $t_i$, we can immediately write the form of the
expansion: % 式 (10.39)
\setcounter{equation}{38}
\begin{equation}\label{eq:bzongo}
\varphi\cdot t_i=\sum_j
(z-z_0)^{m_j-m_i}f_i{}^j\,t_j.
\end{equation}
 For $\varphi$ to be holomorphic at
$y>y_0$, the functions $(z-z_0)^{m_j-m_i}f_i{}^j$ must all be
regular at $z=z_0$.  If we order the $m_i$ in the usual way
$m_1\geq m_2\geq\dots\geq m_N$, then the number of conditions on
the Taylor series coefficients of the $f_i{}^j$ near $z=z_0$ is
$\sum_{i<j}(m_i-m_j)$.

This is the same as the dimension of the space of Hecke
modifications of bundles as determined in eqn. \eqref{eq:ydro}.  This
strongly suggests that, in sharp contrast to the case of Hecke
modifications of bundles, the Hecke modifications of a Higgs
bundle $(E,\varphi)$ are generically a finite set.  In fact, this
is the case if $\varphi(z_0)$ is regular and semisimple (which
means for $G=U(N)$ that it is diagonalizable with distinct
eigenvalues). The most systematic description of the result is to
say that $\varphi$ determines a vector field on the variety $\bar
{\EUY}({}^L\negthinspace w,z_0)$ that parametrizes Hecke
modifications of $E$ of type $^L\negthinspace w$ at the point
$z_0$; the zeroes of this vector field are the possible Hecke
modifications of the pair $(E,\varphi)$. (The vector field is not
quite canonically determined, but its zeroes are.)

Deferring a detailed explanation of this to \cite{ref26}, we will
briefly describe the simple example with $G=U(2)$,
$^L\negthinspace w=(1,0)$. In this example, as we learned in
section 第~\ref{sec:spaceheck}~节, if we take the bundle $E_-$ to be trivial,
then $E_+$ can be characterized by saying that a section of $E_+$
holomorphic at $z=0$ has the form % 式 (10.40)
\setcounter{equation}{39}
\begin{equation}\label{eq:hxonxno}
s_+=(z-z_0)^{-1}
u+(h_1(z),h_2(z)),
\end{equation}
  where $u$ is some nonzero element of
$\Bbb{C}^2$ and $h_1$ and $h_2$ are holomorphic near $z=z_0$.  Up
to scaling by $\Bbb{C}^*$, $u$ parametrizes the space of possible
Hecke modifications of $E$, which is $\Bbb{CP}^1$ in this example.
$\varphi$ is holomorphic at $z=z_0$  for $y>y_0$ if and only if
the product $\varphi\cdot s_+$ is of the same form, or in other
words if and only if $\varphi u$ is a multiple of $u$. In the
regular semisimple case, this means that $u$ must be one of the
two eigenvectors of $\varphi$.

Going back to the extended Bogomolny equations \eqref{eq:BPSextsimple}, an
immediate consequence of the relation
$\partial_y\phi_z=-i[\phi_z,\phi_0]$ is that $\phi_z$ is
everywhere zero if it zero for some $y$.  We assumed this in our
analysis of the product of 't Hooft operators in section 第~\ref{sec:opprod}~节.
A further consequence of the same equation is that the
characteristic polynomial of $\varphi$ is independent of $y$, even
in crossing an 't Hooft operator.  As a result, Hecke
transformations of Higgs bundles map each fiber of the Hitchin
fibration to itself.

Now consider a brane wrapped on a fiber of the Hitchin fibration,
with a flat Chan-Paton bundle. These are the branes that according
to $S$-duality must be magnetic eigenbranes. From what we have
said in the last paragraph, the support of such a brane is
invariant under Hecke transformations.  This observation is a step
toward proving directly, without relying on $S$-duality, that such
branes are magnetic eigenbranes, as will be shown in detail
elsewhere \cite{ref26}.

\par\medskip{\it Hecke Correspondence For Higgs Bundles}

Just as in our discussion of the ordinary Bogomolny equations, we
can consider the possible Hecke modifications of a given Higgs
bundle $(E,\varphi)$ -- which has been our point of view so far --
or we can describe the initial and final states more symmetrically
and describe a Hecke correspondence ${\EUQ}$ on the space of Higgs
bundles. We define ${\EUQ}$ by solving the extended Bogomolny
equations on $W=I\times C$ with suitable boundary conditions at
the two ends and in the presence of specified 't Hooft operators.
We specify both $A$ and $\varphi$ (but not $\phi_0$) at both ends
of $W$, and divide only by gauge transformations that are trivial
at each end.

Let $\EUW_0$ be the space of all pairs $(A,\varphi)$ consisting of
a gauge field and Higgs field on $C$ satisfying $\bar D\varphi=0$.
Then $\EUW_0$ is a gauge theory version of the ``stack'' ${\rm
Higgs}_G$ of all Higgs bundles, just as the space of all
connections is a gauge theory version of ${\rm Bun}_G$, the stack
of all $G$-bundles.  By restricting a solution of the extended
Bogomolny equations to its boundary values, ${\EUQ}$ has a pair of
maps to $\EUW_0$, or in other words, a map to $\EUW_0\times
\EUW_0$. This is the Hecke correspondence for Higgs bundles.

${\EUQ}$ is very complicated, in general, with a variety of
components of different dimension.  However, if we limit ourselves
to a good region of $\EUW_0$ (where each Higgs bundle has only
finitely many possible Hecke modifications of the chosen type up
to a gauge transformation), then ${\EUQ}$ can be regarded as a
submanifold of $\EUW_0\times \EUW_0$. In this good region,
moreover, we can divide by gauge transformations and  ${\EUQ}$
becomes a middle-dimensional submanifold of $\MH\times\MH$.

't Hooft operators preserve supersymmetry of type $(B,A,A)$, as we
explained at the end of section 第~\ref{sec:lineop}~节.  This means that it must
be possible to interpret ${\EUQ}$ as a brane of type $(B,A,A)$ in
$\MH\times \MH$.  In particular, it must be holomorphic in complex
structure $I$ and Lagrangian\footnote{Above the middle dimension,
there are $A$-branes that are not Lagrangian \cite{ref80}. But in the
good region, ${\EUQ}$ is middle-dimensional and must be
Lagrangian.} with respect to complex structures $J$ and $K$ or
equivalently with respect to the holomorphic symplectic form
$\Omega_I=\omega_J+i\omega_K$ defined in eqn. \eqref{eq:ninto}.  More
exactly, if $\pi_i:{\EUQ}\to\MH\times\MH$, $i=1,2$, are the two
projections, then ${\EUQ}$ must be Lagrangian with respect to the
holomorphic symplectic form
$\Omega=\pi_1^*(\Omega_I)-\pi_2^*(\Omega_I)$ on $\MH\times \MH$.
The minus sign here ensures that the diagonal in $\MH\times \MH$,
which is the trivial Hecke correspondence with no 't Hooft
operators, is Lagrangian.

Letting $C_+$ and $C_-$ denote the two ends of $W=I\times C$, we
can write the restriction of $\Omega$ to ${\EUQ}$ as
% 式 (10.41)
\setcounter{equation}{40}
\begin{equation}\label{eq:symplevolume}
{1\over\pi}\left(\int_{C_+}
-\int_{C_-}\right)|d^2 z|\,\Tr\left(\delta\phi_z\wedge \delta
A_\bz\right)={1\over \pi}\int_Idy\int_C|d^2z|\,
\partial_y \Tr\left(\delta\phi_z\wedge\delta A_\bz\right).
\end{equation}

To compute the restriction of this form to $\EUQ$, we can assume
that the variations $\delta\phi_z$ and $\delta A_\bz$ satisfy the
linearization of \eqref{eq:BPSextsimple} . The linearizations of the first
three equations in \eqref{eq:BPSextsimple}\ read
% 式 (10.42)
\setcounter{equation}{41}
\begin{equation}\label{eq:BPSlin}
\begin{aligned}D_\bz\delta\phi_z&=-[\delta
A_\bz,\phi_z],\\
\partial_y\delta A_\bz&=-iD_\bz\delta\phi_0 - i [\delta
A_\bz,\phi_0]\\
\partial_y\delta\phi_z&=-i[\delta\phi_z,\phi_0]-i[\phi_z,\delta\phi_0]
.\\\end{aligned}
\end{equation}
 Using these equations, the Jacobi identity, and integration
by parts, one can verify that \eqref{eq:symplevolume}\ vanishes, so that
${\EUQ}$ is complex-Lagrangian, as expected.
